% ---------------------------------------------------------------------------
% §5 Experimental Design — boş olan Section 5'in tam metni.
% Not: config numaraları YENİ ardışık şemaya göre (eski C7 = yeni C6).
% Vinyet ve soru metinleri Appendix H'ye (appendix_experiment_materials.tex)
% referans verir; \ref etiketlerini kendi appendix düzenine göre güncelle.
% ---------------------------------------------------------------------------

\section{Experimental Design}
\label{sec:experiment}

The experiment asks how the political structure surrounding a protest event
shapes public reactions to it. We manipulate a single factor with two
levels, which we label \emph{political structure} following the
opportunity--threat framework of the political process tradition: a
\textbf{threat} condition, in which the protest occurs during an active
security episode, and an \textbf{opportunity} condition, in which the same
protest occurs during an ongoing peace and negotiation process. Each
condition is operationalized as a short news-style vignette. The threat
vignette reports a Ministry of National Defense statement on intensified
armed clashes in northern Syria, including casualties and heightened
security measures. The opportunity vignette reports ongoing talks between
the government and Kurdish political representatives, widely attended Newroz
celebrations, and Newroz messages issued by the President and party leaders.
Both vignettes close with an identical description of the focal event: a
march in Ankara at which a group voices the demand for Kurdish-language
education. The full Turkish texts of both vignettes, together with English
translations, are reported in Appendix~\ref{app:materials}.

\paragraph{Outcomes.} Two survey items follow each vignette and are given
equal analytical weight. \emph{Movement legitimacy} asks respondents to rate
agreement with the statement that the group is waging a legitimate rights
struggle, on a five-point Likert scale. \emph{Behavioral intent} asks
whether the respondent would personally attend such a march, as a binary
choice. The two items capture complementary faces of public response, an
evaluative judgment about the movement and a costly behavioral disposition,
and mirror the response structures on which the inference protocols were
calibrated: an ordered attitudinal item and a rare binary participation
item. Original item wordings appear in Appendix~\ref{app:materials}.

\paragraph{Inference protocols.} Each outcome is administered under the
protocol calibrated for its response format in Section~4. The Likert
legitimacy item uses the C6 persona configuration with verbalized sampling
at $T=0.8$; the binary intent item uses C6 with direct answering at $T=0$.
No information beyond the persona profile, the vignette, and the survey item
enters the prompt. The prompt templates are identical to the calibration
stage except for the inserted vignette, and are reproduced in
Appendix~\ref{app:materials}.

\paragraph{Assignment and estimation.} The design is within-persona: all
2{,}615 twins respond under both political structures, in independent prompt
turns with no shared conversational state, so that responses in one
condition cannot contaminate the other. Condition-specific respondent-level
seeds control the stochastic sampling step. The primary estimand is the
within-persona paired contrast between the threat and opportunity
conditions, aggregated over the synthetic population; we report paired
$t$-tests with standardized effect sizes ($d_z$) at the population level and
within subgroups. Because each twin serves as its own control, the design
removes all between-respondent heterogeneity from the contrast, which is
precisely the counterfactual that observational protest research cannot
recover: the same public observed under two political structures.

\paragraph{Scope.} Consistent with the calibration results, the experiment
is interpreted at the level of aggregate and subgroup response
distributions. The protocols were validated for distributional resemblance,
not for individual-level prediction, and the experimental claims inherit
that scope.
